\documentclass[11pt]{article}
\usepackage{icmlko}
\begin{document}

\title{채택 셋을 다시 재 봤다 --- 둘은 무리 안에서도 벌고 하나는 어디서도
안 번다}
\author{월드모델 실험실}
\date{노트 472 · 2026년 8월 3일}
\maketitle

\begin{abstract}
노트 470 이 \emph{무리 안 판}을 만들고 노트 471 이 \emph{무리평균만}을
만들었다. 노트 471 은 후자를 가드로 쓰자고 적었는데 \textbf{그건 안
된다} --- 무리평균만 rho 는 \emph{모델과 무관}하므로 ``얹은 몫''$=$
판$-$상수 이고 변화량이 판과 똑같다. 갈라내는 것은 \textbf{무리 안 판}
뿐이다.

그래서 이 사이클의 채택 셋을 그 자로 다시 쟀다(무리 열쇠가 있는 일곱
도메인 · 씨앗 여덟 짝지어).

\begin{center}
\begin{tabular}{lrrl}
\hline
채택 & 판 차 & 무리 안 차 & \\
\hline
계절 (노트 447) & $+0.0040$ & $\mathbf{+0.0044}$ & 같이 벌었다 \\
특기 칸 (노트 449) & $+0.0049$ & $\mathbf{+0.0048}$ & 같이 벌었다 \\
깊이 $12$ (노트 457) & $-0.0016$ & $\mathbf{-0.0023}$ & 둘 다 잃는다 \\
\hline
\end{tabular}
\end{center}

\textbf{무리 사이만 산 채택은 없다} --- 등록한 문장은 통과다. 계절과
특기 칸은 무리 \emph{안}에서 판에서만큼, 혹은 조금 더 번다. 구간도 넷 다
$0$ 을 안 문다.
\end{abstract}

\begin{figure}[t]
\centering
\includegraphics[width=\textwidth]{figs/f1.pdf}
\caption{\textbf{가} --- 채택 셋을 두 자로. \textbf{나} --- 깊이 $12$ 를
지금 있는 일곱 자로 다 재면 초록 둘(수집 전체 범위의 집 밖 짝)만 양수다.}
\end{figure}

\section{깊이 12 만 다르게 생겼다}

깊이 $12$ 가 판을 잃는 것은 새 소식이 아니다 --- 노트 457 이 그것을
\emph{알고} 채택했다(집 밖 두 짝을 사고 판을 조금 잃는 거래, 환율
$1.95$). 새 소식은 \textbf{무리 안 손해가 판 손해보다 크다}는 것이다
($-0.0023$ 대 $-0.0016$).

지금까지 잰 것을 다 모으면 이렇다.

\begin{center}
\begin{tabular}{lr}
\hline
깊이 $6 \to 12$ 가 준 것 & \\
\hline
판 (열한 도메인) & $-0.0032$ \\
판 (일곱 도메인) & $-0.0016$ \\
\textbf{무리 안 판} & $\mathbf{-0.0023}$ \\
KR 만화 · 수집 전체 & $+0.0064$ \\
\textbf{KR 만화 · 갈래 안} & $\mathbf{-0.0090}$ \\
비게임 앱 · 수집 전체 & $+0.0061$ \\
비게임 앱 · 갈래 안 & $+0.0045$ \\
\hline
\end{tabular}
\end{center}

\textbf{일곱 자 중 양수는 둘뿐이고, 그 둘이 채택 근거였다.} 나머지 다섯은
음수이거나(넷) 작다.

\section{그래도 안 되돌린다}

\textbf{되돌리지 않는다.} 이유를 적어 둔다.

\begin{itemize}
\item \textbf{규약의 자는 판(도메인 전체)과 집 밖(수집 전체)이다.} 그
자로 깊이 $12$ 는 노트 426 거울의 조건 넷을 다 채웠고, 노트 469 가
\textbf{그 두 자의 범위가 서로 맞다}는 것을 코드로 확인했다.
\item \textbf{무리 안 판은 어제 만든 자다}(노트 470). 어제 만든 자로
그제 한 채택을 되돌리면 \emph{자를 만들 때마다 과거가 흔들린다.} 노트 470
자신이 ``판은 안 바꾼다''고 적은 것과 같은 이유다.
\item \textbf{노트 469 의 조항} --- 되돌리려면 근거가 기대는 사실을 먼저
확인해야 하는데, \emph{어느 범위가 배포인가}는 아직 사람에게 못 물었다.
\end{itemize}

\begin{quote}
\textbf{대신 미리 적는다} --- 배포가 ``비슷한 것들 사이에서 고른다''로
정해지면 \textbf{깊이 $12$ 를 되돌린다}(이전 상태 이름: ``깊이 $6$ 인
F18''). 그때는 무리 안 판과 갈래 안 전이가 판정 자가 되고, 깊이 $12$ 는
그 둘에서 각각 $-0.0023$ 과 $-0.0090$ 이다.
\end{quote}

\section{계절과 특기 칸은 깨끗하다}

두 채택은 판과 무리 안 판을 \emph{같은 크기로} 올린다. 특히 계절은
무리 안에서 조금 더 번다($+0.0044$ 대 $+0.0040$) --- 연중 위치는 갈래를
맞히는 데 안 쓰이고 갈래 \emph{안}에서 작품을 가르는 데 쓰인다는 뜻이다.
노트 447 이 ``분기 애니가 서로 갉아먹는다''고 읽은 것과 맞는다.

\textbf{이 사이클의 채택 중 둘은 어느 자로 봐도 산다.}

\section*{남은 것}

\begin{itemize}
\item \textbf{어느 범위가 배포인가} --- 이제 답에 따라 무엇이 바뀌는지
까지 적어 뒀으니, 물을 준비가 끝났다.
\item \textbf{무리 열쇠 없는 네 도메인}(도서 · 시장팝업 · 아이돌 · 팝업)
은 이 자로 못 잰다. 판 무게 $12\%$ 다.
\item \textbf{lr · it · K 도 같은 자로 다시} --- 셋 다 판으로만 정했다.
\end{itemize}

\end{document}
